% Generated by scripts/build_reports.py; edit reviewed Markdown inputs.
% Compile with XeLaTeX.
% Options for packages loaded elsewhere
\PassOptionsToPackage{unicode}{hyperref}
\PassOptionsToPackage{hyphens}{url}
\PassOptionsToPackage{dvipsnames,svgnames,x11names}{xcolor}
\documentclass[
  11pt,
]{article}
\usepackage{xcolor}
\usepackage[margin=1in]{geometry}
\usepackage{amsmath,amssymb}
\setcounter{secnumdepth}{-\maxdimen} % remove section numbering
\usepackage{iftex}
\ifPDFTeX
  \usepackage[T1]{fontenc}
  \usepackage[utf8]{inputenc}
  \usepackage{textcomp} % provide euro and other symbols
\else % if luatex or xetex
  \usepackage{unicode-math} % this also loads fontspec
  \defaultfontfeatures{Scale=MatchLowercase}
  \defaultfontfeatures[\rmfamily]{Ligatures=TeX,Scale=1}
\fi
\usepackage{lmodern}
\ifPDFTeX\else
  % xetex/luatex font selection
\fi
% Use upquote if available, for straight quotes in verbatim environments
\IfFileExists{upquote.sty}{\usepackage{upquote}}{}
\IfFileExists{microtype.sty}{% use microtype if available
  \usepackage[]{microtype}
  \UseMicrotypeSet[protrusion]{basicmath} % disable protrusion for tt fonts
}{}
\makeatletter
\@ifundefined{KOMAClassName}{% if non-KOMA class
  \IfFileExists{parskip.sty}{%
    \usepackage{parskip}
  }{% else
    \setlength{\parindent}{0pt}
    \setlength{\parskip}{6pt plus 2pt minus 1pt}}
}{% if KOMA class
  \KOMAoptions{parskip=half}}
\makeatother
\usepackage{longtable,booktabs,array}
\usepackage{caption}
\captionsetup[table]{skip=6pt}
\newcounter{none} % for unnumbered tables
\usepackage{calc} % for calculating minipage widths
% Correct order of tables after \paragraph or \subparagraph
\usepackage{etoolbox}
\makeatletter
\patchcmd\longtable{\par}{\if@noskipsec\mbox{}\fi\par}{}{}
\makeatother
% Allow footnotes in longtable head/foot
\IfFileExists{footnotehyper.sty}{\usepackage{footnotehyper}}{\usepackage{footnote}}
\makesavenoteenv{longtable}
\setlength{\emergencystretch}{3em} % prevent overfull lines
\providecommand{\tightlist}{%
  \setlength{\itemsep}{0pt}\setlength{\parskip}{0pt}}
\usepackage{amsmath,amssymb}
\usepackage{xurl}
\usepackage{fancyhdr}
\pagestyle{fancy}
\fancyhf{}
\fancyhead[L]{\small Pure Thought | Revised research specifications}
\fancyfoot[L]{\small 2026-09-17}
\fancyfoot[R]{\small\thepage}
\setlength{\headheight}{15pt}
\setlength{\emergencystretch}{3em}
\allowdisplaybreaks
\usepackage{bookmark}
\IfFileExists{xurl.sty}{\usepackage{xurl}}{} % add URL line breaks if available
\urlstyle{same}
% fallback for those not using the hyperref driver hyperxmp:
\makeatletter
\@ifundefined{xmpquote}{\newcommand{\xmpquote}[1]{#1}}{}
\makeatother
\hypersetup{
  pdftitle={03. Celestial amplitudes: distributional crossing and soft-limit consistency},
  pdfauthor={Pure Thought research portfolio},
  colorlinks=true,
  linkcolor={blue},
  filecolor={Maroon},
  citecolor={Blue},
  urlcolor={blue},
  pdfcreator={LaTeX via pandoc}}

\title{03. Celestial amplitudes: distributional crossing and soft-limit
consistency}
\author{Pure Thought research portfolio}
\date{Revised 17 September 2026}

\begin{document}
\maketitle
\thispagestyle{empty}
\clearpage

\section{03. Celestial amplitudes: distributional crossing and
soft-limit consistency}

\textbf{Revised specification --- 2026-09-11.}
\href{https://github.com/rick-stevens-ai/Pure-Thought/blob/main/reviews/03-Celestial-CFT-Bootstrap.md}{Detailed
critique} \ensuremath{\cdot}
\href{https://github.com/rick-stevens-ai/Pure-Thought/blob/main/archive/original-PRDs/03-Celestial-CFT-Bootstrap.md}{Original
description} \ensuremath{\cdot}
\href{https://github.com/rick-stevens-ai/Pure-Thought/blob/main/reviews/EVIDENCE-STANDARD.md}{Evidence
standard}

\subsection{Plain-language guide}

\subsubsection{The problem in everyday terms}

Instead of describing a collision by particle energies and momenta,
imagine recording the directions in which particles enter or leave on a
sphere surrounding the experiment. Celestial methods reorganize
scattering data into quantities associated with that sphere. The
ambition is a new language for gravity. Before using that language to
discover new theories, this problem checks that it faithfully preserves
familiar facts about a simple graviton collision.

\subsubsection{Key terms}

\begin{itemize}
\item
  \textbf{Celestial sphere} --- The sphere of directions seen from a
  point, used here to label asymptotic particle directions rather than a
  physical shell.
\item
  \textbf{Graviton} --- The hypothetical quantum excitation of the
  gravitational field used in perturbative quantum-gravity calculations.
\item
  \textbf{Helicity} --- The component of a particle's spin along its
  direction of motion.
\item
  \textbf{Mellin transform} --- A mathematical transformation that
  reorganizes dependence on energy into dependence on scaling behavior.
\item
  \textbf{Conformal weight \ensuremath{\Delta}} --- A label describing
  scaling behavior in the celestial description; it can be complex.
\item
  \textbf{Distribution} --- A generalized function, such as a delta
  function, understood through integration against test functions rather
  than ordinary pointwise values.
\item
  \textbf{Crossing} --- The relationship between different scattering
  processes obtained by reinterpreting an incoming particle as an
  outgoing antiparticle, with appropriate continuation.
\item
  \textbf{Soft limit} --- The limit in which one particle's energy
  approaches zero; a soft theorem controls the resulting leading
  behavior.
\item
  \textbf{Regulator} --- A temporary modification that makes a singular
  expression well defined, accompanied by a rule for removing or
  interpreting it.
\end{itemize}

\subsubsection{Why this matters}

A successful reformulation could expose symmetries or constraints that
are difficult to see in ordinary scattering variables. That would give
theorists another way to investigate gravity. But a change of
mathematical language is useful only if singularities, conservation laws
and limiting cases survive the translation correctly.

\subsubsection{What progress would mean}

A verified example would supply a dependable foundation for later
celestial calculations. It would establish that specific identities
translate correctly, rather than establish a complete two-dimensional
theory of the universe.

\subsection{Core question}

Can a regulated celestial transform of a specified four-graviton tree
amplitude reproduce Lorentz covariance, crossing, and the leading
conformally soft residue in one consistent convention?

\subsection{Scope and assumptions}

Fix four-dimensional Einstein gravity at tree level, a nonzero
four-graviton helicity amplitude, incoming/outgoing signs
\ensuremath{\epsilon}\ensuremath{{}_{i}}, and
p\ensuremath{{}_{i}}=\ensuremath{\epsilon}\ensuremath{{}_{i}}\ensuremath{\omega}\ensuremath{{}_{i}}q(z\ensuremath{{}_{i}},\ensuremath{\bar{z}}\ensuremath{{}_{i}}).
Declare metric, helicity, normalization, and analytic-continuation
conventions. The target is an identity pipeline for amplitudes, not a
classification of celestial CFTs.

Use
\ensuremath{\Delta}\ensuremath{{}_{i}}=1+i\ensuremath{\lambda}\ensuremath{{}_{i}}
on the principal series where appropriate and
h\ensuremath{{}_{i}}=(\ensuremath{\Delta}\ensuremath{{}_{i}}+J\ensuremath{{}_{i}})/2,
\ensuremath{\bar{h}}\ensuremath{{}_{i}}=(\ensuremath{\Delta}\ensuremath{{}_{i}}\ensuremath{-}J\ensuremath{{}_{i}})/2.
Identify explicitly which continuations away from that contour are used
for soft residues.

\subsection{Mathematical target}

Define the transformed distribution by

\[\widetilde{\mathcal A}=\prod_i\int_0^\infty d\omega _i\,\omega _i^{\Delta _i-1}\;\delta ^{(4)}\!\left(\sum_i \epsilon _i\omega _iq_i\right)\mathcal M_4.\]

State a regulator or test-function space making the integral meaningful,
and specify regulator removal or meromorphic continuation. From a
leading momentum-space term proportional to
\ensuremath{\omega}\_s\ensuremath{{}^{-1}}, derive the Mellin pole at
\ensuremath{\Delta}\_s=1 in this convention. Subleading residues require
their own derivation. Crossing must transform helicities, signs, and
support, not just permute ordinary Euclidean functions.

\subsection{Required outputs and proof obligations}

Produce exact momentum-space expressions, regulated transform formulas,
support conditions, and symbolic covariance/crossing/soft identities. If
a transform cannot be defined under the chosen assumptions, return a
precise obstruction at that step. A positive bootstrap is an extension
only after deriving an applicable inner product and positivity
statement.

\subsection{Validation and rejection controls}

Check the original amplitude's factorization first. Test Mellin formulas
against analytic one-variable examples in their convergence strips.
Verify soft residues before and after transformation under stated
limit-interchange hypotheses. Use smeared distributions for numerical
checks; never evaluate a delta distribution as an ordinary function.

\subsection{Milestones and research extension}

Baseline: one four-point transform with explicit support and checked
leading soft residue. Strong: a second helicity/process example and a
verified crossing relation. Research extension: derive a justified
positive constraint in a restricted subsector before attempting an SDP
island.

Milestones are acceptance gates, not calendar promises. Before a run,
freeze the concrete instance/domain, tolerance, compute budget, and
permitted dependencies. A bounded search that fails to find a result
must preserve an unresolved outcome.

\subsection{Starting instruction}

Work on the baseline above. First make every input and convention
explicit, then validate the reference case. Build the checker alongside
the calculation. Label numerical evidence, finite-model certificates,
and continuum theorems separately; return the strongest justified result
and identify the exact unmet proof obligation.

\subsection{Primary references and their role}

\begin{itemize}
\tightlist
\item
  \href{https://www.pure.ed.ac.uk/ws/files/121747399/Adamo_2019_Class._Quantum_Grav._36_205018_1_.pdf}{Adamo
  et al., Celestial amplitudes and conformal soft theorems} ---
  Conformally soft graviton operators at \ensuremath{\Delta}=1 and
  \ensuremath{\Delta}=0.
\end{itemize}

References support definitions or limitations; the scope choices and
proposed milestones are editorial recommendations, not claims that these
results have already been obtained.

\subsection{Displayed formulation and claim boundary}

\href{https://github.com/rick-stevens-ai/Pure-Thought/blob/main/PRDs/03-Celestial-CFT-Bootstrap.md}{Canonical
specification}.

\[p_i=\epsilon_i\omega_i q_i,\quad
\Delta_i=1+i\lambda_i,\quad
h_i=\frac{\Delta_i+J_i}{2},\quad
\bar h_i=\frac{\Delta_i-J_i}{2}.\]

\[\widetilde{\mathcal A}=
\left(\prod_{i=1}^4\int_0^\infty d\omega_i\,\omega_i^{\Delta_i-1}\right)
\delta^{(4)}\!\left(\sum_i\epsilon_i\omega_iq_i\right)\mathcal M_4.\]

The expression is a regulated distribution. A leading soft term
proportional to \(\omega_s^{-1}\) yields the local integral
\(\int_0^\Lambda\omega_s^{\Delta_s-2}d\omega_s=\Lambda^{\Delta_s-1}/(\Delta_s-1)\)
before continuation, exposing the pole at \(\Delta_s=1\). Support,
helicity, crossing and regulator removal remain proof obligations;
ordinary Euclidean OPE positivity is not assumed.

\subsection{Detailed review and rationale}

\href{https://github.com/rick-stevens-ai/Pure-Thought/blob/main/PRDs/03-Celestial-CFT-Bootstrap.md}{Revised
problem} \ensuremath{\cdot}
\href{https://github.com/rick-stevens-ai/Pure-Thought/blob/main/archive/original-PRDs/03-Celestial-CFT-Bootstrap.md}{Preserved
original} \ensuremath{\cdot}
\href{https://github.com/rick-stevens-ai/Pure-Thought/blob/f854d4e12b55c1f5b5c2a964b75f2d23496d981e/PRDs/03-Celestial-CFT-Bootstrap.md}{Pinned
upstream source}

Line references below refer to the preserved original at commit
\texttt{f854d4e12b55}, not the rewritten file. ``Blocking'' means the
original issue can invalidate the scientific target or claimed
certificate; ``major'' means it materially changes scope, inference, or
acceptance criteria. These are critiques of specifications, not results
of running the embedded research-code sketches.

\subsubsection{Detailed findings}

\begin{enumerate}
\def\labelenumi{\arabic{enumi}.}
\item
  \textbf{Blocking --- missing distributional specification (36--45,
  94--113).} A Mellin transform of a stripped amplitude without momentum
  conservation and convergence rules is incomplete. The numerical
  quadrature sketch does not define the multidimensional distribution
  that is actually needed.
\item
  \textbf{Blocking --- imported positivity (18--23, 57--75).}
  Principal-series celestial data cannot simply inherit the positive
  OPE-coefficient squares of a unitary Euclidean two-dimensional CFT. An
  SDP needs a proved positivity structure, not an asserted positive
  norm.
\item
  \textbf{Major --- soft poles (60--64).} The generic
  \ensuremath{\Delta}\ensuremath{\to}0 double-pole formula does not
  reproduce the leading graviton soft pole in the stated Mellin
  convention. Deriving the pole from the energy power prevents
  convention mistakes.
\item
  \textbf{Major --- scope (15--28, 79--86).} ``Consistent space'' and
  ``islands'' assume tools whose foundational prerequisites are
  themselves research questions. A distributionally correct identity
  pipeline is a meaningful prerequisite and preserves the original
  direction.
\end{enumerate}

\subsubsection{What the revision preserves and changes}

The central motivation is retained, but the executable question is now:
\textbf{Can a regulated celestial transform of a specified four-graviton
tree amplitude reproduce Lorentz covariance, crossing, and the leading
conformally soft residue in one consistent convention?} The revised
scope chooses a concrete starting model; those choices are proposed
research-design decisions, not facts inferred from the original. The
baseline can succeed without resolving the full open research program.

The replacement also removes the original implementation sketches from
the active specification. They remain in the archive for comparison, but
have not been made executable or certified. New implementations should
derive their equations from the revised definitions rather than copy
unchecked prototypes.

\subsubsection{Acceptance criteria to use instead}

Baseline: one four-point transform with explicit support and checked
leading soft residue. Strong: a second helicity/process example and a
verified crossing relation. Research extension: derive a justified
positive constraint in a restricted subsector before attempting an SDP
island.

\subsubsection{Sources supporting the corrections}

\begin{itemize}
\tightlist
\item
  \href{https://www.pure.ed.ac.uk/ws/files/121747399/Adamo_2019_Class._Quantum_Grav._36_205018_1_.pdf}{Adamo
  et al., Celestial amplitudes and conformal soft theorems} ---
  Conformally soft graviton operators at \ensuremath{\Delta}=1 and
  \ensuremath{\Delta}=0.
\end{itemize}

\section{Evidence and execution standard}

This standard applies to the revised specifications. It is a proposed
contract for future research runs; this editorial revision does not
supply the scientific certificates themselves.

\subsection{What ``pure thought'' means here}

Use explicitly stated mathematical models and deductions without fitting
conclusions to experimental outcomes. Synthetic examples, exact model
inputs, analytic references, numerical search, and established software
are allowed. Physical constants, electronic-structure approximations,
constitutive laws, and empirical parameter sets must be identified when
used. A model-derived result is conditional on that model.

The original ban on legacy software conflicts with its proposed reliance
on NumPy, SciPy, CVXPY, nauty, and other libraries. Reimplementing every
optimizer and algebra system is not necessary for scientific
independence. Prefer a small independently checkable certificate and a
documented trusted base. A fresh implementation can be a project
objective, but must not be confused with a mathematical proof or novelty
claim.

\subsection{Four evidence levels}

{\def\LTcaptype{none} % do not increment counter
\begin{longtable}[]{@{}
  >{\raggedright\arraybackslash}p{(\linewidth - 4\tabcolsep) * \real{0.3333}}
  >{\raggedright\arraybackslash}p{(\linewidth - 4\tabcolsep) * \real{0.3333}}
  >{\raggedright\arraybackslash}p{(\linewidth - 4\tabcolsep) * \real{0.3333}}@{}}
\toprule\noalign{}
\begin{minipage}[b]{\linewidth}\raggedright
Label
\end{minipage} & \begin{minipage}[b]{\linewidth}\raggedright
What it establishes
\end{minipage} & \begin{minipage}[b]{\linewidth}\raggedright
What it does not establish
\end{minipage} \\
\midrule\noalign{}
\endhead
\bottomrule\noalign{}
\endlastfoot
Numerical evidence & A reproducible computation at stated precision,
sample size, and discretization & Exact truth, continuum coverage, or a
universal theorem \\
Finite-model certificate & An exact/validated statement about the
specified finite model or finite relaxation & An omitted tail,
thermodynamic limit, or physical realization \\
Conditional mathematical theorem & A proof over the stated domain with
every hypothesis and error controlled & Truth outside those assumptions
or a demonstrated physical system \\
Formal proof & A checked theorem in a named proof assistant, with
imported axioms/dependencies listed & Correct physical modeling merely
by virtue of formalization \\
\end{longtable}
}

Formalization is an optional extra deliverable unless the specific
project makes it mandatory. High precision, JSON output, solver status
and small residuals do not by themselves confer any of these proof
levels.

\subsection{Claim directions}

For minimization over a larger feasible set, the optimum is a lower
bound on the original minimum. A physical feasible point supplies an
upper bound. For maximization the directions reverse. Track feasibility
of the primal/dual witnesses separately from their objective values.

Infeasibility of an inner ansatz excludes only that ansatz.
Infeasibility of an outer relaxation can exclude the original problem if
the map into the relaxation is proved and all approximations preserve
necessity. Passing necessary tests means ``not excluded''; existence
needs a constructive sufficient argument.

A smallest observed example is not a minimality theorem. A complete
catalog must define the universe, equivalence relation, boundaries, and
proof that pruning misses nothing. A finite-size numerical trend does
not prove an asymptotic limit.

\subsection{Minimum run manifest}

Before calculation, record the problem ID and revision hash, exact
inputs and conventions, quantified domain, equivalence relation, target
observable, accepted evidence level, tolerances, and resource limits.
Record dependencies and precision/rounding mode. Where randomization is
used, preserve seeds, trial counts and raw sufficient statistics.
Resource limits are planning decisions, not evidence that a mathematical
solution cannot exist.

An output bundle should include:

\begin{itemize}
\tightlist
\item
  The statement being claimed, including all assumptions and
  quantifiers.
\item
  Construction or proof object, plus an independently runnable verifier
  and its trust assumptions.
\item
  Numerical/error ledger separating roundoff, discretization,
  truncation, optimization, sampling and modeling error where relevant.
\item
  Positive reference cases and negative controls that the verifier must
  reject.
\item
  Coverage: included domain, omitted domain, unresolved regions and
  termination reason.
\end{itemize}

These are requirements for future implementations, not a promise of
polynomial-size certificates for every problem.

\subsection{Checking specific certificate types}

\textbf{LP/SDP/SOS:} verify equality constraints, objective direction,
and cone feasibility using exact arithmetic or outward-rounded
intervals. Rationalizing a numerical solution can break PSD constraints.
Correct residuals with justified norm bounds; do not silently project
and reuse the old objective. Distinguish a proof of a polynomial
approximation from a proof for the original function.

\textbf{Interval numerics:} state function enclosures, domains and
outward rounding. A small root box needs existence and, if claimed,
uniqueness. Cover the complement for a global enumeration. Include
singularities and tails rather than discarding them numerically.

\textbf{SAT/SMT:} a solver's UNSAT label is not a certificate. Preserve
a proof supported by a checker and the proof that the encoding, symmetry
reduction and blocked solutions match the mathematical universe. DRAT
applies to an appropriate Boolean CNF proof, not automatically to every
real-arithmetic or SDP claim.

\textbf{Statistics:} specify the estimand, dependence assumptions,
confidence procedure and bias. Zero observed failures does not prove
zero failure probability. MCMC error bars based on estimated
autocorrelation are not generally rigorous mixing guarantees.

\subsection{Milestones and novelty}

Baseline means a working, independently validated result on a concrete
reference problem. Strong means a broader or sharper result with its
claimed evidence level. Research extension means a possible novel result
whose feasibility and significance need review. A negative theorem can
be valuable; an unsuccessful search can also be useful if its scope is
documented, but the two are not interchangeable.

Do not promise ``publication-quality'' in a fixed number of months.
Before claiming novelty, conduct a focused current literature comparison
against the exact theorem/model/metric. The sources in this revision
check key definitions and limitations; they are not an exhaustive
novelty audit across all 30 fields.

\end{document}
